%==========================================================================
\section{Comparison with previous prescriptions}
\label{app:choices}
%==========================================================================
% Must be \input BEFORE appendix_params.tex: it is the first appendix cited
% in the text (Sect. 2), so it becomes Appendix A and app:params Appendix B.
% Figures: same files as before (2_model/t_cs_vgrav.pdf, produced by
% t_cs_vgrav.py; 2_model/t_align_oloc.pdf, produced by t_align_oloc.py).

The model of Sect.~\ref{sec:model} follows \citet{Motta2017} for the precession frequency and the wave-propagation time, and \citet{FoucartLai2014} for the alignment rate, with the three differences listed in Table~\ref{tab:choices}.
The first, the number of cycles required by condition (ii), is discussed in Sect.~\ref{sec2:timescales}.
In this appendix we derive the statements made in Sect.~\ref{sec2:timescales} on the other two, the sound speed (Sect.~\ref{appM:cs}) and the local precession rate (Sect.~\ref{appM:omloc}), and quantify their effect on the timescales.
The notation is that of Sect.~\ref{sec:model}.

\begin{table}
\caption{Prescriptions adopted in this work and in the reference papers.}
\label{tab:choices}
\centering
\begin{tabular}{lccc}
\hline\hline
 & FL14 & M18 & This work \\
\hline
Condition (ii)          & --                          & $t_\textup{align}>2/\nu_0$   & $t_\textup{align}>1/\nu_0$ \\
$c_s$                   & $H\,\Omega_z$               & $H\,\Omega_\textup{K}$       & $H\,\Omega_z$ \\
$\Omega_\textup{loc}$   & $\mathcal{Z}\Omega_\varphi$ & $\mathcal{Z}\Omega_\varphi$  & $\Omega_\textup{LT}$ \\
$\omega_\textup{p}$     & $\langle\mathcal{Z}\Omega_\varphi\rangle$ & $\langle\Omega_\textup{LT}\rangle$ & $\langle\Omega_\textup{LT}\rangle$ \\
Eq.~\eqref{eq2:BC}      & holds                       & violated                     & holds \\
\hline
\end{tabular}
\tablefoot{
FL14: \citet{FoucartLai2014}; M18: \citet{Motta2017}.
$\langle\cdot\rangle$ is the mean over $[R_\textup{in},R_\textup{out}]$ with the weight $\tilde w=\Sigma R^3\Omega_\varphi$, so that $\langle\Omega_\textup{LT}\rangle=\Omega_\textup{solid}$ (Eq.~\ref{eq2:nusolid}).
For FL14, $\omega_\textup{p}$ is the value required by Eq.~\eqref{eq2:BC}.
}
\end{table}
% check the short labels FL14 / M18 against the years in biblio.bib

%--------------------------------------------------------------------------
\subsection{Sound speed}
\label{appM:cs}

\citet{FoucartLai2014} define the sound speed that sets both the propagation speed of bending waves and the alignment rate of Eq.~\eqref{eq2:talign} through the vertically integrated pressure, Eq.~\eqref{eq2:csdef}, so that $c_s=H\,\Omega_z$.
\citet{Motta2017} instead set $H=c_s/\Omega_\textup{K}$ with $c_s(R_g)=c\,(H/R)_0$ (their Sect.~4.1), where $\Omega_\textup{K}\doteq t_g^{-1}r^{-3/2}$.
%This normalisation requires $R\,\Omega_\textup{K}=c\,r^{-1/2}$, i.e.\ the Newtonian $\Omega_\textup{K}=(c/R_g)\,r^{-3/2}$, although in their Eqs.~(2) and (7) the same symbol denotes the relativistic orbital frequency.
The two definitions of $c_s$ coincide in Newtonian gravity but not around a spinning hole.

\paragraph{Sign of $\Omega_z-\Omega_\textup{K}$.}
Let $u\doteq a\,r^{-3/2}$.
Dividing numerator and denominator of Eq.~\eqref{eq2:epiphi} by $r^{3/2}$ gives $\Omega_\varphi=t_g^{-1}r^{-3/2}/(1+u)$, i.e.\ $\Omega_\varphi/\Omega_\textup{K}=1/(1+u)$, and Eq.~\eqref{eq2:epiz} then gives
\begin{equation}
    \label{eqM:OzOK}
    \left(\frac{\Omega_z}{\Omega_\textup{K}}\right)^2=\frac{1-4u+3a^2r^{-2}}{(1+u)^2}.
\end{equation}
Since $(1+u)^2>0$, $\Omega_z<\Omega_\textup{K}$ if and only if $1-4u+3a^2r^{-2}<1+2u+u^2$, i.e., with $u^2=a^2r^{-3}$, $a^2r^{-2}\left(3-r^{-1}\right)<6\,a\,r^{-3/2}$.
For $a>0$, dividing by $a\,r^{-3/2}>0$ gives $a\,r^{-1/2}(3-r^{-1})<6$, which holds because $0<a\leq1$ and $r\geq r_\textup{ISCO}\geq1$ bound the left-hand side by three.
For $a<0$ the left-hand side is positive and the right-hand side negative, so the inequality is reversed.
Hence, at every $r\geq r_\textup{ISCO}$,
\begin{equation}
    \label{eqM:OzOKsign}
    \Omega_z<\Omega_\textup{K}\quad(a>0),
    \qquad
    \Omega_z>\Omega_\textup{K}\quad(a<0).
\end{equation}
Expanding Eq.~\eqref{eqM:OzOK} for $r\gg1$ gives $\Omega_z/\Omega_\textup{K}\simeq1-3a\,r^{-3/2}$.

\paragraph{Effect on the timescales.}
At fixed $H$, by Eq.~\eqref{eqM:OzOKsign} the Newtonian prescription overestimates $c_s$ at every radius for prograde spins and underestimates it for retrograde ones.
Since $t_\textup{wave}$ is an integral of $1/c_s$ (Eq.~\ref{eq2:twave_int}), and $t_\textup{align}$ increases with $c_s$ at every radius (Eq.~\ref{eq2:talign}, where $\mathcal{G}$ and $\Sigma$ do not depend on $c_s$), the prescription $c_s=H\Omega_\textup{K}$ makes both conditions (i) and (ii) of Sect.~\ref{sec2:timescales} easier to satisfy for $a>0$ and harder for $a<0$.
Figure~\ref{fig:cs_choice} shows the effect on both timescales.
At $r_\textup{out}=3\,r_\textup{ISCO}$, Eq.~\eqref{eq2:csdef} lengthens $t_\textup{wave}$ by factors of 1.95 and 1.16 for $a=0.998$ and $0.8$, and shortens $t_\textup{align}$ by factors of 3.8 and 1.4; for $a=-1$ it shortens $t_\textup{wave}$ by $4\%$ and lengthens $t_\textup{align}$ by $10\%$.
The difference is largest for compact flows around rapidly spinning holes, and becomes negligible at large radii.

\begin{figure}
    \centering
    \includegraphics[width=\hsize]{2_model/t_cs_vgrav.pdf}
    \caption{
    Wave-propagation time (top) and alignment time (bottom) per unit mass as functions of $r_\textup{out}$, for $c_s=H\Omega_z$ (solid; this work, Eq.~\ref{eq2:csdef}) and $c_s=H\Omega_\textup{K}$ (dashed; \citealt{Motta2017}), with the same $H$ and the spins in the legend ($a=1$ is actually $a=0.998$, so that $t_\textup{wave}$ is not divergent).
    Both use $\Omega_\textup{loc}=\Omega_\textup{LT}$, $h_\textup{ISCO}=0.1$, $\alpha=0.01$, $p=3/5$ and $q=3/2$.
    }
    \label{fig:cs_choice}
\end{figure}

\paragraph{Thickness profile.}
In the formulation of \citet{Motta2017}, $c_s=c\,h_0\,r^{-q}$ and $H=c_s/\Omega_\textup{K}$.
With $1/\Omega_\textup{K}=t_g\,r^{3/2}$ and $c\,t_g=R_g$ this gives $H=h_0\,r^{3/2-q}R_g$.
The choice $q=3/2$ is therefore equivalent to a thickness independent of radius.
We retain this property, $H=h_\textup{ISCO}\,r_\textup{ISCO}\,R_g$ (Eq.~\ref{eq2:HR}), and obtain the sound speed from Eq.~\eqref{eq2:csdef}, which gives Eq.~\eqref{eq2:cs}.
Since $H$ is normalised at each spin's own ISCO, a retrograde flow is thicker than a prograde one with the same $|a|$ (Sect.~\ref{sec2:timescales}).
At fixed $r$ the ratio of their sound speeds is larger than that of their thicknesses, because by Eq.~\eqref{eqM:OzOKsign} $\Omega_z(-|a|)>\Omega_\textup{K}>\Omega_z(|a|)$.

%--------------------------------------------------------------------------
\subsection{Local precession rate}
\label{appM:omloc}

\citet{FoucartLai2014} work with a quasi-Keplerian potential and take $\Omega_\textup{loc}=\mathcal{Z}\Omega_\varphi$ in Eq.~\eqref{eq2:G}, with (their Eq.~16)
\begin{equation}
    \label{eqM:Zdef}
    \mathcal{Z}\doteq\frac{\Omega_\varphi^2-\Omega_z^2}{2\Omega_\varphi^2}
    =\frac12\left[1-\left(\frac{g_z}{g_\varphi}\right)^{2}\right].
\end{equation}

\paragraph{The closure of \citet{Motta2017}.}
\citet{Motta2017} adopt Eq.~\eqref{eq2:talign} with $\Omega_\textup{loc}=\mathcal{Z}\Omega_\varphi$, but take $\omega_\textup{p}=\Omega_\textup{solid}$ from Eq.~\eqref{eq2:nusolid}, i.e.\ the $\tilde w$-mean of $\Omega_\textup{LT}$ (their Eqs.~2 and 7).
This violates Eq.~\eqref{eq2:BC}.
Using $\Omega_z=\Omega_\varphi-\Omega_\textup{LT}$ (Eq.~\ref{eq2:omlt}), $\Omega_\varphi^2-\Omega_z^2=(\Omega_\varphi-\Omega_z)(\Omega_\varphi+\Omega_z)=\Omega_\textup{LT}\,(2\Omega_\varphi-\Omega_\textup{LT})$, hence, exactly,
\begin{equation}
    \label{eqM:LTminusZ}
    \Omega_\textup{LT}-\mathcal{Z}\Omega_\varphi=\frac{\Omega_\textup{LT}^2}{2\Omega_\varphi}.
\end{equation}
Splitting $\omega_\textup{p}-\mathcal{Z}\Omega_\varphi=(\Omega_\textup{solid}-\Omega_\textup{LT})+(\Omega_\textup{LT}-\mathcal{Z}\Omega_\varphi)$, their stress is
\begin{equation}
    \label{eqM:GM}
    \mathcal{G}_\textup{M}(R)=\mathcal{G}_\textup{LT}(R)+E(R),
    \qquad
    E(R)\doteq\frac12\int_{R_\textup{in}}^{R}\Sigma x^3\,\Omega_\textup{LT}^2\,\textup{d}x ,
\end{equation}
where $\mathcal{G}_\textup{LT}$ is Eq.~\eqref{eq2:G} with $\Omega_\textup{loc}=\Omega_\textup{LT}$ and $\omega_\textup{p}=\Omega_\textup{solid}$, and we used $\tilde w\,\Omega_\textup{LT}^2/(2\Omega_\varphi)=\Sigma x^3\Omega_\textup{LT}^2/2$.
Since $\Omega_\textup{solid}$ is the $\tilde w$-mean of $\Omega_\textup{LT}$, $\mathcal{G}_\textup{LT}(R_\textup{out})=\Omega_\textup{solid}\int\tilde w\,\textup{d}x-\int\tilde w\,\Omega_\textup{LT}\,\textup{d}x=0$, and therefore
\begin{equation}
    \label{eqM:GMout}
    \mathcal{G}_\textup{M}(R_\textup{out})=E(R_\textup{out})>0\qquad(a\neq0).
\end{equation}
The stress then does not vanish at the outer edge, contrary to the boundary condition from which Eq.~\eqref{eq2:talign} is derived.

\paragraph{Motivation of our choice.}
We take $\Omega_\textup{loc}=\Omega_\textup{LT}$, together with $\Omega_\textup{solid}$ from Eq.~\eqref{eq2:nusolid} (Eq.~\ref{eq2:omloc}), for three reasons.

First, Eq.~\eqref{eq2:BC} then holds by construction for every $\Sigma$, $a$ and $r_\textup{out}$, and the scheme is formally that of \citet{FoucartLai2014} with $\mathcal{Z}\Omega_\varphi$ replaced by $\Omega_\textup{LT}$.

Second, the linear theory cannot distinguish between the two consistent choices.
Writing $\Omega_z=\Omega_\varphi(1-\eta)$, Eq.~\eqref{eq2:omlt} gives $\Omega_\textup{LT}=\eta\,\Omega_\varphi$, and Eq.~\eqref{eqM:Zdef} gives $\mathcal{Z}\Omega_\varphi=\Omega_\varphi(\eta-\eta^2/2)$, since $\Omega_\varphi^2-\Omega_z^2=\Omega_\varphi^2(2\eta-\eta^2)$.
The two choices thus coincide at first order in the deviation from Keplerian motion, with relative difference $\eta/2=\Omega_\textup{LT}/(2\Omega_\varphi)$.
The bending-wave equations require $|\Omega_z^2-\Omega_\varphi^2|\lesssim(H/R)\,\Omega_\varphi^2$ \citep{FoucartLai2014}, i.e.\ $2|\mathcal{Z}|\lesssim H/R$.
Since $\Omega_z\geq0$ implies $\eta\leq1$, we have $|2-\eta|\geq1$ and hence $2|\mathcal{Z}|=|\eta|\,|2-\eta|\geq|\eta|$.
Wherever the equations hold, the two choices therefore differ by $|\eta|/2\leq|\mathcal{Z}|\lesssim(H/R)/2$, which is of the same order as the terms in $H/R$ neglected in deriving the equations.

Third, the remaining ambiguity is fixed by the single-particle limit, which the model is required to contain (Sect.~\ref{sec2:pif}).
Since $\Omega_\textup{loc}$ is continuous on $[R_\textup{in},R_\textup{out}]$ and $\tilde w>0$, the $\tilde w$-mean in Eq.~\eqref{eq2:BC} lies between the minimum and the maximum of $\Omega_\textup{loc}$ on that interval, and therefore tends to $\Omega_\textup{loc}(R_\textup{in})$ as $R_\textup{out}\to R_\textup{in}$.
With $\Omega_\textup{loc}=\Omega_\textup{LT}$, a narrow ring therefore precesses at the exact nodal frequency of the geodesics it is made of.
With $\Omega_\textup{loc}=\mathcal{Z}\Omega_\varphi$ it would precess at $\Omega_\textup{LT}(R_\textup{in})\left[1-\Omega_\textup{LT}/(2\Omega_\varphi)\right]$ (Eq.~\ref{eqM:LTminusZ}).
This underestimates the nodal frequency by $39\%$ at the ISCO of an $a=0.998$ hole, where $\Omega_\textup{LT}/\Omega_\varphi\simeq0.77$ (Eqs.~\ref{eq2:epiphi}--\ref{eq2:omlt}).
Our choice thus uses the same rate in the observable frequency, Eq.~\eqref{eq2:nusolid}, and in the stress, and recovers the exact test-particle limit in both.

\paragraph{Validity close to the ISCO.}
The relative difference $\Omega_\textup{LT}/(2\Omega_\varphi)$ is $\simeq a\,r^{-3/2}$ for $r\gg1$ and vanishes for $a\to0$, but it becomes of order unity close to the ISCO of a fast-spinning hole.
There, however, the validity condition above is violated anyway: $2\mathcal{Z}=-\varepsilon$ (Sect.~\ref{sec2:signs}) equals $\simeq0.57$ at $r_\textup{ISCO}$ for $a=0.9$ and $\simeq0.95$ for $a=0.998$, both well above $h_\textup{ISCO}$.
The analogous condition on $\kappa$ fails close to the ISCO for every spin, since $\kappa(r_\textup{ISCO})=0$.
However, $\kappa$ enters the first-order solution of \citet{FoucartLai2014} only through the periastron term in the warp profile and amplitude (their Eqs.~28 and 34), not through $\mathcal{G}$ or $\gamma$ (their Eqs.~31 and 35).
In the innermost region of a fast-spinning hole, the choice of $\Omega_\textup{loc}$ is therefore a systematic uncertainty of the linear theory that cannot be settled within it.

\paragraph{Size and sign of the difference.}
Both closures share the denominator of Eq.~\eqref{eq2:talign}, $\Sigma$ and $c_s$, so they differ only through $\mathcal{G}^2$, with, by Eq.~\eqref{eqM:GM}, $\mathcal{G}_\textup{M}^2-\mathcal{G}_\textup{LT}^2=E\,(E+2\mathcal{G}_\textup{LT})$.
By Eq.~\eqref{eq2:signLT}, $\mathcal{G}_\textup{LT}=s\,\mathcal{G}_\textup{abs}$ with $s\doteq\textup{sign}(a)$ (Sect.~\ref{sec2:retro}), while $E\geq0$ is even in $s$.
Moreover, $\mathcal{G}_\textup{abs}\leq0$ wherever $|\Omega_\textup{LT}|$ decreases with $r$, which we verified numerically on $r_\textup{ISCO}\leq r\leq500$ for all the spins considered.
Indeed, $\textup{d}\mathcal{G}_\textup{abs}/\textup{d}R=\tilde w\,(|\Omega_\textup{solid}|-|\Omega_\textup{LT}|)$ changes sign once, from negative to positive, and $\mathcal{G}_\textup{abs}$ vanishes at both edges.
For $a<0$, then, $E+2\mathcal{G}_\textup{LT}=E+2|\mathcal{G}_\textup{abs}|\geq0$ at every radius: the closure of \citet{Motta2017} overestimates $\gamma$, and our $t_\textup{align}$ is longer at every $r_\textup{out}$.
For $a>0$ the two terms compete, and the sign of the difference depends on $r_\textup{out}$.
Figure~\ref{fig:omloc_choice} shows the result.
For $a=-1$ our $t_\textup{align}$ is longer by a factor of 1.11 at $r_\textup{out}=3\,r_\textup{ISCO}$ and by $4\%$ at $r_\textup{out}=500$.
For $a>0$ it is shorter, except in very compact flows around rapidly spinning holes ($r_\textup{out}\lesssim1.5\,r_\textup{ISCO}$ for $a=0.998$, $\lesssim1.1\,r_\textup{ISCO}$ for $a=0.8$).
The difference reaches a factor of about five for $a=0.998$ at $r_\textup{out}\simeq2.3\,r_\textup{ISCO}$, and stays below $5\%$ for $|a|\leq0.2$ and $r_\textup{out}\geq3\,r_\textup{ISCO}$.

The term $E$ also breaks the symmetry of Sect.~\ref{sec2:retro}: since $\mathcal{G}_\textup{M}=s\,\mathcal{G}_\textup{abs}+E$ with $E$ even in $s$, $\mathcal{G}_\textup{M}^2$ depends on the sign of the torque, whereas $\mathcal{G}_\textup{LT}^2=\mathcal{G}_\textup{abs}^2$ does not.

\begin{figure}
    \centering
    \includegraphics[width=\hsize]{2_model/t_align_oloc.pdf}
    \caption{
    Alignment time per unit mass as a function of $r_\textup{out}$ for $\Omega_\textup{loc}=\Omega_\textup{LT}$ (solid; this work) and for $\Omega_\textup{loc}=\mathcal{Z}\Omega_\varphi$ with $\Omega_\textup{solid}$ from Eq.~\eqref{eq2:nusolid} (dashed; \citealt{Motta2017}), for the spins in the legend, with $c_s=H\Omega_\textup{K}$ (top) and $c_s=H\Omega_z$ (bottom).
    Both panels use $h_\textup{ISCO}=0.1$, $\alpha=0.01$, $p=3/5$ and $q=3/2$.
    For $a=1$, $\Omega_z\to0$ at the ISCO, so that $\mathcal{Z}\Omega_\varphi\to\Omega_\varphi/2$ while $\Omega_\textup{LT}\to\Omega_\varphi$: with $\Omega_\textup{loc}=\mathcal{Z}\Omega_\varphi$ the difference $\Omega_\textup{solid}-\Omega_\textup{loc}$ near the ISCO is comparable to $\Omega_\varphi$ instead of a residual that averages to zero, the internal torque does not vanish at $R_\textup{out}$, and $t_\textup{align}$ is shortened (dashed below solid for $r_\textup{out}\lesssim2$).
    With $c_s=H\Omega_z$ the sound speed also vanishes at the ISCO, $c_s^2\propto(R-R_\textup{ISCO})^2$, which cancels the $(r_\textup{out}-r_\textup{ISCO})^{-2}$ divergence of compact flows: the dashed curve then tends to a finite value as $r_\textup{out}\to r_\textup{ISCO}$, whereas the solid one still diverges.
    }
    \label{fig:omloc_choice}
\end{figure}